% Part: lambda-calculus
% Chapter: introduction
% Section: lambda-computable

\documentclass[../../../include/open-logic-section]{subfiles}

\begin{document}

\olfileid[fr]{lam}{int}{cmp}
\olsection{Les fonctions \texorpdfstring{$\lambd$}{λ}-définissables sont calculables}

\begin{thm}
\ollabel{thm:lambda-computable}
Toute fonction partielle~$f$ définie par un lambda-terme
est partielle calculable.
\end{thm}

\begin{proof}
Supposons que la fonction~$f$ soit définie par un
lambda-terme~$X$. Décrivons une procédure informelle
pour calculer~$f$. Étant donnés $m_0$, \dots,
$m_{n-1}$, écrivons le terme
$X\,\num{m_0}\,\dots\,\num{m_{n-1}}$.
Construisons un arbre : plaçons d'abord toutes les
réductions en une étape du terme initial, puis, au
niveau suivant, toutes les réductions en une étape
de ces termes, c'est-à-dire les réductions en deux
étapes du terme initial, et poursuivons ainsi.
Parcourons l'arbre niveau par niveau, en commençant par le terme initial. Si nous atteignons
un numéral, à $\alpha$-équivalence près, renvoyons l'entier qu'il représente.
Si aucun numéral n'est jamais atteint, la recherche
continue indéfiniment : $f$ n'est pas définie pour
ces arguments.

En invoquant la thèse de Church, nous pouvons conclure
que cette fonction est partielle calculable. Une
meilleure démonstration donnerait une description
récursive de cette recherche. On pourrait, par exemple,
définir une suite de fonctions et de relations primitives
récursives, « $\fn{IsASubterm}$ »,
« $\fn{Substitute}$ »,
« $\fn{ReducesToInOneStep}$ »,
« $\fn{ReductionSequence}$ »,
« $\fn{Numeral}$ », etc. Pour calculer
$f(m_0, \dots, m_{n-1})$, la procédure récursive
partielle recherche alors une suite de réductions
en une étape qui part de
$X\,\num{m_0}\,\dots\,\num{m_{n-1}}$ et aboutit
à un numéral, puis renvoie l'entier correspondant.
Les détails sont longs et fastidieux, mais ne présentent
pas d'autre difficulté.
\end{proof}

\end{document}
